\documentclass[12pt,a4paper,oneside]{report}

% -----------------------------------------------------------------------------
% Transit-fitting analysis report for the TESS exoplanet pipeline
% -----------------------------------------------------------------------------
\usepackage[a4paper,margin=1in]{geometry}
\usepackage[T1]{fontenc}
\usepackage[utf8]{inputenc}
\usepackage{lmodern}
\usepackage{microtype}
\usepackage{amsmath,amssymb,mathtools}
\usepackage{booktabs,longtable,array}
\usepackage{graphicx}
\usepackage{xcolor}
\usepackage{enumitem}
\usepackage{listings}
\usepackage{hyperref}
\usepackage{cleveref}
\usepackage{fancyhdr}
\usepackage{caption}

\hypersetup{colorlinks=true,linkcolor=blue!55!black,citecolor=green!45!black,urlcolor=blue!55!black}
\newcommand{\SI}[2]{\ensuremath{#1\,\mathrm{#2}}}
\newcommand{\um}[1]{#1}
\setlist{itemsep=2pt,topsep=4pt}
\setlength{\parindent}{0pt}
\setlength{\parskip}{5pt}
\pagestyle{fancy}
\fancyhf{}
\fancyhead[L]{TESS transit-fitting pipeline}
\fancyhead[R]{\thepage}
\fancyfoot[C]{\small Reproducible technical report}

\newcommand{\Msun}{\ensuremath{M_{\odot}}}
\newcommand{\Rsun}{\ensuremath{R_{\odot}}}
\newcommand{\Rearth}{\ensuremath{R_{\oplus}}}
\newcommand{\Rstar}{\ensuremath{R_{\star}}}
\newcommand{\Mstar}{\ensuremath{M_{\star}}}
\newcommand{\Teff}{\ensuremath{T_{\mathrm{eff}}}}
\newcommand{\rhoStar}{\ensuremath{\rho_{\star}}}
\newcommand{\Rp}{\ensuremath{R_{\mathrm{p}}}}
\newcommand{\Porb}{\ensuremath{P}}
\newcommand{\Tfourteen}{\ensuremath{T_{14}}}
\newcommand{\tzero}{\ensuremath{t_{0}}}
\newcommand{\dd}{\,\mathrm{d}}
\newcommand{\code}[1]{\texttt{\small #1}}

\lstset{basicstyle=\ttfamily\small,breaklines=true,frame=single,columns=fullflexible,showstringspaces=false}

\title{\textbf{Transit-Fitting Analysis Pipeline}\\[0.5em]
\large A detailed technical report for the TESS exoplanet-modelling workflow}
\author{TESS Exoplanet Pipeline}\date{15 July 2026}

\begin{document}
\maketitle
\pagenumbering{roman}
\tableofcontents
\clearpage
\pagenumbering{arabic}

\chapter*{Executive summary}
\addcontentsline{toc}{chapter}{Executive summary}
This report documents the complete transit-fitting workflow implemented in the
\code{tess\_pipeline} package. The workflow starts with a TIC identifier or a
local TESS FITS light curve and ends with posterior distributions for transit
parameters, derived planetary properties, convergence diagnostics, plots, and
machine-readable results. The primary inference path is a joint PyMC--
\code{exoplanet} model sampled with the No-U-Turn Sampler (NUTS). A Gaussian
process (GP) accounts for time-correlated stellar and instrumental variability,
while the transit signal is generated from a Keplerian orbit and a quadratic
limb-darkened light-curve model.

The analysis is deliberately staged. Archive information and period searches
provide initial timing information; TESS photometry supplies the transit shape;
Gaia/TIC/SIMBAD information supplies stellar constraints; and the final
Bayesian model propagates uncertainty jointly rather than treating a point
estimate as exact. Unless explicitly overridden, the documented defaults are a
TESS SPOC two-minute cadence, TLS period detection, a four-chain NUTS run with
1500 posterior draws and 1000 tuning draws per chain, target acceptance 0.95,
and a stochastically-driven harmonic oscillator (SHO) GP kernel.

\textbf{Important scope.} This document describes the implemented model and its
assumptions, including implementation choices that should not be interpreted as
universal astronomical standards. In particular, the current code uses a
single fitted flux offset for the Bayesian light curve even though the prior
helper is written to support multiple sectors, and the reported BIC parameter
count follows the project's fixed convention. These details are recorded for
reproducibility.

\chapter{Scientific objective and data model}
\section{Objective}
The objective is to determine whether periodic flux decreases in TESS
photometry are consistent with one or more transiting planets and, conditional
on that interpretation, estimate
\begin{itemize}
  \item orbital period \Porb and reference transit time \tzero,
  \item radius ratio \(r=\Rp/\Rstar\), impact parameter \(b\), and duration \Tfourteen,
  \item stellar density \rhoStar and orbital scale \(a/\Rstar\),
  \item physical planet radius, semi-major axis, and equilibrium temperature,
  \item uncertainty intervals, residual structure, and sampling reliability.
\end{itemize}
The observable is a time series \(\{t_i,F_i,\sigma_i\}_{i=1}^{N}\), where
\(t_i\) is normally BTJD, \(F_i\) is normalized relative flux, and \(\sigma_i\)
is the reported per-cadence uncertainty. The likelihood uses the original
uncertainty when available and adds an inferred white-noise jitter term.

\section{Entry points and execution modes}
A complete run is equivalent to:
\begin{lstlisting}
analysis = TESSAnalysis("TIC 307210830", inference=True)
analysis.resolve_target()
analysis.lookup_archive_period()
analysis.load_lightcurves()
analysis.preprocess()
analysis.search_period()
analysis.query_gaia()
analysis.characterize_star()
analysis.query_sdss()
analysis.fit_transit()
analysis.derive_planet_parameters()
analysis.check_convergence()
analysis.generate_figures()
analysis.save("output/")
\end{lstlisting}
The same stages are executed by \code{analysis.run()}. The inference backend
can be \code{exoplanet} (the primary Bayesian path), \code{batman\_only} (a
fast analytic fit without MCMC), or \code{phoebe} (an advanced binary/star-
planet modelling path). This report concentrates on the \code{exoplanet}
backend, while documenting the alternatives where relevant.

\section{Default configuration}
\begin{longtable}{p{0.36\textwidth}p{0.18\textwidth}p{0.36\textwidth}}
\caption{Principal defaults used by the pipeline.}\label{tab:defaults}\\
\toprule
Parameter & Default & Meaning\\
\midrule
\code{author} & SPOC & TESS light-curve product author\\
\code{cadence} & 120 s & Two-minute short cadence\\
\code{search\_method} & TLS & Transit Least Squares\\
\code{period\_min,max} & 0.5, 100 d & Broad period-search interval\\
\code{max\_planets} & 1 & Maximum blind-search candidates\\
\code{sigma\_clip\_lower,upper} & 20, 5 & Asymmetric outlier thresholds\\
\code{flatten\_window\_length} & 401 & Savitzky--Golay window in cadences\\
\code{flatten\_polyorder} & 3 & Polynomial order for flattening\\
\code{flatten\_break\_tolerance} & 5 & Gap tolerance in cadences\\
\code{chains} & 4 & Independent MCMC chains\\
\code{draws} & 1500 & Retained draws per chain\\
\code{tune} & 1000 & Adaptation/tuning draws per chain\\
\code{target\_accept} & 0.95 & NUTS target acceptance probability\\
\code{gp\_kernel} & SHO & Correlated-noise covariance kernel\\
\bottomrule
\end{longtable}
All values can be changed through explicit constructor arguments, a TOML
configuration file, or \code{TESS\_PIPELINE\_*} environment variables. Explicit
arguments have highest priority.

\chapter{End-to-end pipeline}
\section{Target resolution}
The target stage normalizes a TIC identifier supplied as an integer, ``TIC
123'', or ``TIC123''. In download mode it resolves the TIC to a target record
and sky coordinates (right ascension and declination). In FITS mode, the
identifier and coordinates may be read from FITS headers. The resolved target
is retained in the results metadata so that every catalogue query and output
can be traced to the same object.

\section{Archive information and period initialization}
The pipeline optionally queries the NASA Exoplanet Archive. A user-supplied
period override takes precedence over the archive. If archive periods exist,
each is refined in a narrow interval around the literature value. If no archive
period is available, the search covers the configured interval.

The archive is used as initialization and provenance, not as a replacement for
photometric fitting. The final period has a continuous prior with a one-percent
width around each supplied candidate unless the model is explicitly configured
to fix it. The epoch prior is centered on the archive/search epoch with a
standard deviation of 
um{0.1} day.

\section{TESS light-curve acquisition}
The default product is a SPOC light curve obtained through \code{lightkurve}
from MAST. The user may instead provide one or more local FITS files. Sectors
are recorded after loading, and the raw sectors are retained for diagnostic
plots. The expected photometric product is PDCSAP-like relative flux, although
the implementation accepts a compatible \code{LightCurve} object.

The following assumptions apply at this stage:
\begin{enumerate}
  \item time values are comparable across sectors after lightkurve loading;
  \item quality flags have been handled by the selected lightkurve product and
        its quality bitmask;
  \item the target is sufficiently isolated that the extracted aperture flux is
        a useful proxy for the target flux; and
  \item flux uncertainties are either supplied by the light curve or replaced
        by a conservative default scale of \(10^{-3}\) in the Bayesian routine.
\end{enumerate}

\section{Preprocessing and normalization}
Preprocessing is applied per sector before stitching:
\begin{enumerate}
  \item remove NaN cadences;
  \item construct a transit mask if \Porb and \tzero are available;
  \item sigma-clip only non-transit points, protecting candidate transits;
  \item stitch the cleaned sectors;
  \item flatten the stitched curve with a Savitzky--Golay filter;
  \item remove NaNs and perform a final outlier pass.
\end{enumerate}
For a time \(t\), phase is defined by
\begin{equation}
  \phi(t)=\left[\frac{t-\tzero}{\Porb}\right]_{\mathrm{wrapped}},
  \qquad -0.5\leq \phi < 0.5.
\end{equation}
The default protected region is \(|\phi|<0.15/\Porb\), equivalent to a
window of \SI{0.15}{day} in time around each predicted transit. This is a
protection mask, not a measurement of the final transit duration.

Outliers are removed with lower and upper thresholds of 20 and 5 standard
deviations, respectively. The asymmetric choice is intentional: it is more
aggressive for positive excursions than for negative transit-like excursions.
The default flattening window has 401 two-minute cadences, approximately
\SI{13.4}{hour}. The transit mask is passed to the flattening operation so that
the first-pass trend is not trained on the signal being measured. The resulting
curve is normalized near unit flux; residual correlated variability is left for
the GP in the final fit.

\section{Period and transit detection}
\subsection{Transit Least Squares and Box Least Squares}
TLS searches a grid of periods and transit durations using limb-darkened transit
shapes. It is the preferred detector for small, non-box-shaped signals. BLS is
available as an alternative or cross-check and uses a box-shaped template. The
canonical detection record contains \Porb, \tzero, \Tfourteen, depth, method,
SDE or SNR, false-alarm probability (FAP), confidence, and a search-based
existence probability.

For a TLS detection with signal detection efficiency \(\mathrm{SDE}\), the
implemented trial-corrected approximation is
\begin{equation}
  \mathrm{FAP}=1-\left(1-e^{-\mathrm{SDE}}\right)^{n_{\mathrm{trials}}},
  \qquad P_{\mathrm{search}}=1-\mathrm{FAP}.
\end{equation}
The number of trials is the number of TLS periods when available, otherwise
1000. Confidence labels are ``High'' for FAP \(\leq 0.05\), ``Moderate'' for
\(0.05<\mathrm{FAP}\leq0.50\), and ``Low'' above 0.50. These are screening
labels and are not a formal Bayesian model probability.

\subsection{Multiple candidates and aliases}
Blind multi-planet searches are iterative. After a candidate is found, cadences
within \(0.75\Tfourteen\) of each predicted transit are masked and the residual
light curve is searched again. A broad search may downsample a very large curve
to ten-minute cadence. The allowed period maximum is capped when possible so
that at least three transits fit in the time baseline. A coarse SDE/SNR below
4 stops the search; additional candidates must pass a refined SDE/SNR threshold
of 6.

To reduce harmonic confusion, the code refines the initial period near
\(fP\), where \(f\in\{1/3,1/2,1,2,3\}\), and retains the candidate with the
largest refined SDE/SNR. This procedure is a practical alias-resolution rule;
it does not prove that the selected period is dynamically unique.

\section{Stellar characterization}
Stellar information enters the transit model primarily through the stellar
density prior and is also required to convert a radius ratio into a physical
planet radius. Gaia DR3 is queried first for the target position. Missing
quantities are supplemented by a local NASA TOI archive and then by the MAST
TIC catalogue. The characterization module additionally queries VizieR TIC
v8.2 and SIMBAD and merges values with the following parameter-level priority:
\begin{enumerate}
  \item VizieR TIC v8.2;
  \item Gaia DR3;
  \item SIMBAD.
\end{enumerate}
The special priority for metallicity is SIMBAD spectroscopic [Fe/H], followed
by VizieR and Gaia. Provenance, reference, units, and reported uncertainties
are retained in the results.

The stellar quantities are \(\Rstar\), \(\Mstar\), \Teff, \(\log g\), [Fe/H],
parallax, luminosity, and \rhoStar. If radius is unavailable but luminosity and
\Teff are available, it is estimated from the Stefan--Boltzmann relation:
\begin{equation}
  \frac{\Rstar}{\Rsun}=\left(\frac{L_\star}{L_\odot}\right)^{1/2}
  \left(\frac{T_{\odot}}{\Teff}\right)^2.
\end{equation}
As a last resort the implementation uses a main-sequence \Teff scaling,
\(\Rstar/\Rsun=(\Teff/5772\,\mathrm{K})^2\), and assigns a 30-percent radius
uncertainty. If mass is missing, the code uses the approximate scaling
\begin{equation}
  \frac{\Mstar}{\Msun}=
  \left(\frac{\Teff}{5777\,\mathrm{K}}\right)^{1.5}
  \left(\frac{\Rstar}{\Rsun}\right)^{0.1}.
\end{equation}
The density is calculated from mass and radius as
\begin{equation}
  \rhoStar=\rho_\odot\frac{\Mstar/\Msun}{(\Rstar/\Rsun)^3},
  \qquad \rho_\odot\simeq\SI{1.408}{g\,cm^{-3}}.
\end{equation}
Fallback estimates must be flagged in scientific interpretation because their
uncertainty can dominate the planet-radius and orbital-geometry inference.

\chapter{Bayesian transit model}
\section{Observed-flux model}
The primary model represents the normalized flux at time \(t_i\) as
\begin{equation}
  F_{\mathrm{obs}}(t_i)=F_{\mathrm{tr}}(t_i;\boldsymbol{\theta}_{\mathrm{tr}})
  + \mu + g(t_i)+\epsilon_i,
\end{equation}
where \(F_{\mathrm{tr}}\) is the sum of limb-darkened planet transits, \(\mu\)
is a fitted mean-flux offset, \(g(t)\) is correlated GP variability, and
\(\epsilon_i\) is independent Gaussian measurement noise plus jitter. In the
code, the transit model is represented as
\begin{equation}
  F_{\mathrm{tr}}(t)=1+\mu+\sum_{k=1}^{N_{\mathrm{p}}}\Delta F_k(t),
\end{equation}
and the GP prediction is added to obtain \(F_{\mathrm{model}}
=F_{\mathrm{tr}}+g\). The transit light curves \(\Delta F_k\) are negative
flux perturbations generated by \code{exoplanet}.

\section{Orbital geometry}
Each candidate is assigned a Keplerian orbit with period \Porb, transit epoch
\tzero, impact parameter \(b\), stellar density \rhoStar, and radius ratio
\(r=\Rp/\Rstar\). Eccentricity is not sampled explicitly in this implementation;
the default interpretation is a circular transit geometry. Consequently,
transit-derived density should be interpreted as a circular-orbit density unless
an eccentricity model is added.

For a circular orbit, Kepler's third law in stellar units gives
\begin{equation}
  \left(\frac{a}{\Rstar}\right)^3
  =\frac{\rhoStar\,G\,\Porb^2}{3\pi}
  =52.91572\,\rhoStar\,\Porb^2,
\end{equation}
when \rhoStar is in \(\mathrm{g\,cm^{-3}}\) and \Porb is in days. The impact parameter
is related approximately to inclination by \(b=(a/\Rstar)\cos i\) for a
circular orbit. The full orbital geometry is evaluated by
\code{xo.orbits.KeplerianOrbit}, rather than by this approximation alone.

\section{Limb darkening}
The stellar intensity is represented by a quadratic limb-darkening law,
\begin{equation}
  \frac{I(\mu)}{I(1)}=1-u_1(1-\mu)-u_2(1-\mu)^2,
  \qquad \mu=\cos\vartheta.
\end{equation}
Directly sampling \(u_1,u_2\) can produce unphysical combinations. The model
therefore samples the Kipping variables \(q_1,q_2\sim\mathcal{U}(0,1)\) and
transforms them as
\begin{equation}
  u_1=2\sqrt{q_1}\,q_2,
  \qquad
  u_2=\sqrt{q_1}(1-2q_2).
\end{equation}
The same pair is shared by all planets in one fit, because limb darkening is a
property of the host star and not of an individual planet.

\section{Priors}
The prior distributions actually used by the primary model are listed in
\cref{tab:priors}. A uniform prior on a logarithm is a log-uniform prior on the
positive physical parameter.
\begin{longtable}{p{0.24\textwidth}p{0.31\textwidth}p{0.35\textwidth}}
\caption{Bayesian priors in the implemented PyMC model.}\label{tab:priors}\\
\toprule
Quantity & Prior & Interpretation\\
\midrule
\(\Porb{}_{k}\) & \(\mathcal{N}(P_{\mathrm{det},k},(0.01P_{\mathrm{det},k})^2)\) & One-percent period refinement\\
\(\tzero{}_{k}\) & \(\mathcal{N}(t_{0,\mathrm{det},k},0.1^2\,\mathrm{d}^2)\) & Timing refinement\\
\(\log r_k\) & \(\mathcal{U}(\log0.001,\log0.5)\) & Radius-ratio bounds\\
\(b_k\) & \(\mathcal{U}(0,1+r_k)\) & Transit geometry\\
\(q_1,q_2\) & \(\mathcal{U}(0,1)\) & Kipping limb darkening\\
\(\log\Tfourteen{}_{k}\) & \(\mathcal{U}(\log0.01,\log0.5)\) & Duration in days, if fitted\\
\(\log\sigma_{\mathrm{GP}}\) & \(\mathcal{N}(-3,2^2)\) & GP amplitude\\
\(\log\rho_{\mathrm{GP}}\) & \(\mathcal{N}(\log10,0.75^2)\) & GP correlation timescale in days\\
\(\mu\) & \(\mathcal{N}(0,0.01^2)\) & Mean-flux offset\\
\(\log s\) & \(\mathcal{N}(-6,2^2)\) & White-noise jitter scale\\
\bottomrule
\end{longtable}
Here \(s=\exp(\log s)\), \(\sigma_{\mathrm{GP}}=\exp(\log\sigma_{\mathrm{GP}})\),
and \(\rho_{\mathrm{GP}}=\exp(\log\rho_{\mathrm{GP}})\). The period is
technically mutable in the current Bayesian call: \code{period\_fixed=False}
is passed by the inference routine.

\section{Stellar-density treatment and duration parameterization}
The default \code{fit\_duration=True} path samples \Tfourteen directly. This
avoids forcing duration to be determined only by a stellar-density point
estimate. For each planet, the code computes
\begin{align}
  s_k &= \sin\left(\frac{\pi\Tfourteen{}_{k}}{\Porb{}_{k}}\right),\\
  \left(\frac{a}{\Rstar}\right)_k
    &=\left[b_k^2+\frac{(1+r_k)^2-b_k^2}{s_k^2}\right]^{1/2},\\
  \rho_{\star,k} &=\frac{(a/\Rstar)_k^3}{52.91572\,\Porb{}_{k}^2},\\
  \rhoStar &= \frac{1}{N_{\mathrm{p}}}\sum_k\rho_{\star,k}.
\end{align}
A PyMC potential applies a Gaussian stellar-density constraint,
\begin{equation}
  \log p_{\rho}=\log\mathcal{N}(\rhoStar\mid
  \rho_{\star,\mathrm{cat}},\sigma_{\rho,\mathrm{cat}}).
\end{equation}
If no catalogue density is available, the fallback is
\(\rhoStar\sim\operatorname{LogNormal}(\log1.4,1)\). When
\code{fit\_duration=False}, the alternative path samples \rhoStar directly as a
truncated normal around the catalogue value, or as the same broad log-normal
fallback. This distinction is important when comparing analyses.

\section{Gaussian-process noise}
The GP is marginalized jointly with the transit parameters. Its covariance
matrix is added to the diagonal observational variance:
\begin{equation}
  K_{ij}=K_{\mathrm{GP}}(t_i,t_j)+
  \delta_{ij}\left(\sigma_i^2+s^2\right).
\end{equation}
The default celerite2 SHO kernel has quality factor
\(Q=1/\sqrt{2}\), fixed in the code. The alternative is a Matérn
\(\nu=3/2\) kernel. The GP timescale prior is intentionally centered at ten
days with width 0.75 in log space so that the GP is less able to absorb short
transit events. This is a prior regularization assumption, not a guarantee that
stellar variability has a ten-day physical period.

For residual vector \(\boldsymbol{r}=\boldsymbol{F}-\boldsymbol{m}
(\boldsymbol{\theta})\), the marginalized Gaussian-process log likelihood is
\begin{equation}
  \log p(\boldsymbol{F}\mid\boldsymbol{\theta})
  =-\frac{1}{2}\left[\boldsymbol{r}^{\mathsf T}K^{-1}\boldsymbol{r}
  +\log|K|+N\log(2\pi)\right].
\end{equation}
The celerite2 implementation exploits the structured kernel for efficient
factorization and prediction. The posterior also stores a GP prediction,
transit model, full flux model, and residuals for diagnostics.

\section{Likelihood and posterior}
Let \(\boldsymbol{\theta}\) denote all fitted quantities. The posterior is
\begin{equation}
  p(\boldsymbol{\theta}\mid\boldsymbol{F})
  \propto p(\boldsymbol{F}\mid\boldsymbol{\theta})
  p(\boldsymbol{\theta})p_{\rho}(\boldsymbol{\theta}),
\end{equation}
where the density factor is implemented as a PyMC \code{Potential}. The model
is sampled with NUTS, a Hamiltonian Monte Carlo method that adapts a trajectory
length to avoid random-walk behavior. Initialization uses \code{adapt\_diag}
and physically motivated initial values: \(b=0.1\), \(q_1=q_2=0.3\),
\(\log s=-6\), \(\log\sigma_{\mathrm{GP}}=-3\),
\(\log\rho_{\mathrm{GP}}=\log10\), and \(\log r\) initialized from the
square root of the detected depth.

\chapter{Multi-planet model selection}
\section{Candidate-probability gate}
For a second or later detection, the inference stage estimates a rough existence
probability from its SDE/SNR:
\begin{equation}
  p_2=\frac{1}{1+\exp[-0.5(\mathrm{SDE}_2-7)]}.
\end{equation}
A multi-planet MCMC fit is launched only when \(p_2>0.75\). Otherwise, the
candidate is recorded as ``Not Justified'' and the fit is trimmed to the first
planet. This sigmoid is a computational gate, not a calibrated posterior
probability from the light curve.

\section{BIC comparison}
When the gate passes, a one-planet model is sampled first. A final multi-planet
model is then sampled and compared with
\begin{equation}
  \mathrm{BIC}=k\ln N-2\ln\mathcal{L}_{\max},
\end{equation}
where \(N\) is the number of unique time samples and the implementation uses
\(k=7+4N_{\mathrm{p}}\): seven base parameters plus four per planet. The
reported difference is
\begin{equation}
  \Delta\mathrm{BIC}=\mathrm{BIC}_{N}-\mathrm{BIC}_{1},
  \qquad
  p_{N}=\frac{1}{1+\exp(\Delta\mathrm{BIC}/2)}.
\end{equation}
Negative \(\Delta\mathrm{BIC}\) favors the multi-planet model. The mapping to
\(p_N\) is a convenient relative-evidence summary and should not be confused
with a rigorously normalized model posterior. The first candidate is labeled
``Confirmed'' when retained; secondary labels are Very Strong, Strong, Positive,
Weak, or Not Justified according to the size of \(\Delta\mathrm{BIC}\).

\chapter{Posterior-derived physical parameters}
\section{Posterior summaries}
For each planet, samples are flattened over chains and draws. The pipeline
reports the median, standard deviation, and 16th and 84th percentiles for fitted
and derived quantities. The interval \([q_{16},q_{84}]\) is a central
approximately one-standard-deviation credible interval only when the posterior
is close to symmetric; it remains a quantile interval for skewed posteriors.

\section{Transit geometry}
For every posterior sample, the code derives
\begin{equation}
  \frac{a}{\Rstar}=\left(52.91572\,\rhoStar\,\Porb^2\right)^{1/3}.
\end{equation}
The total first-to-fourth-contact duration is computed as
\begin{equation}
  \Tfourteen=\frac{\Porb}{\pi}
  \arcsin\left[\left(
  \frac{(1+r)^2-b^2}{(a/\Rstar)^2-b^2}
  \right)^{1/2}\right],
\end{equation}
then converted from days to hours by multiplying by 24. Arguments are clipped
to the physical interval before evaluating the square root and arcsine.

\section{Planet radius and semi-major axis}
Using the stellar radius in solar units,
\begin{equation}
  \frac{\Rp}{\Rearth}=r\frac{\Rstar}{\Rsun}\frac{\Rsun}{\Rearth}
  =109.203\,r\frac{\Rstar}{\Rsun}.
\end{equation}
The semi-major axis in astronomical units is calculated from \(a/\Rstar\):
\begin{equation}
  \frac{a}{\mathrm{AU}}=
  \frac{(a/\Rstar)(\Rstar/\Rsun)}{215.032}.
\end{equation}
The current implementation propagates stellar-radius uncertainty by drawing
normal samples around \Rstar and \Teff, truncating nonphysical radius and
temperature draws to positive floors. This is Monte Carlo uncertainty
propagation, not a joint stellar posterior sampled in the transit model.

\section{Equilibrium temperature}
The reported equilibrium temperature assumes zero Bond albedo and complete
redistribution of absorbed energy over the planet:
\begin{equation}
  T_{\mathrm{eq}}=\Teff\left(\frac{1}{4(a/\Rstar)^2}\right)^{1/4}.
\end{equation}
This expression uses \(T_\star(a)\approx\Teff/\sqrt{a/\Rstar}\) and the
factor \(1/4\) for full redistribution. It is not a prediction of the actual
surface temperature, atmospheric temperature, or day-side temperature.

\chapter{Validation and diagnostics}
\section{Sampling convergence}
The convergence stage uses ArviZ summaries while excluding large deterministic
vectors such as the full light curve, GP prediction, and model curves. The
reported checks are:
\begin{itemize}
  \item maximum rank-normalized split \(\hat R\), with a threshold of 1.01;
  \item minimum bulk effective sample size (ESS), with a threshold of 400;
  \item total number of divergent NUTS transitions, which must be zero.
\end{itemize}
The run is marked converged only if all available checks pass. A failure should
prompt inspection of trace plots, increasing \code{target\_accept}, extending
tuning, improving parameterization, checking priors, or revisiting the light
curve and stellar constraints. Convergence does not establish that the
astrophysical model is correct; it only indicates that the sampler explored its
specified posterior acceptably.

\section{Photometric and posterior checks}
The generated products include raw and flattened light curves, periodograms,
phase-folded data, Bayesian fit curves, corner plots, trace plots, posterior
predictive curves, residual plots, GP residual autocorrelation, and stacked
individual transits. The following checks are recommended:
\begin{enumerate}
  \item verify that the transit appears in more than one event and more than one
        sector when the baseline permits;
  \item compare the fitted transit depth with the raw and flattened flux;
  \item inspect whether the GP follows out-of-transit variability rather than
        attenuating the transit itself;
  \item inspect residual autocorrelation before and after GP correction;
  \item check that posterior predictive curves cover the observed data without
        leaving coherent phase-dependent residuals;
  \item compare TLS and BLS period evidence and examine harmonics;
  \item verify that the inferred \rhoStar is compatible with catalogue values
        and that fallback stellar estimates have not driven the result.
\end{enumerate}

\section{Known limitations and failure modes}
\begin{itemize}
  \item A transit search is not a validation of planetary origin. Blends,
        eclipsing binaries, stellar activity, dilution, and instrumental systematics
        require additional vetting.
  \item The default orbital interpretation is circular. Eccentric systems can
        yield biased density and duration if eccentricity is not modeled.
  \item The current GP is global in time and uses one fitted mean-flux element
        in the Bayesian call. Sector-specific offsets may be required for strong
        discontinuities between sectors.
  \item The initial flattening is a preprocessing operation. An overly short
        window can suppress or distort long-duration transits; an overly long
        window can leave systematics for the GP.
  \item The catalogue merge is parameter-wise. Correlations among stellar
        radius, mass, temperature, and density are not fully propagated.
  \item The BIC approximation uses a fixed parameter-count convention and a
        maximum sampled likelihood, so it should be used comparatively and not
        as a replacement for full marginal likelihood model comparison.
  \item The archive period and the search-based existence probability encode
        prior information and detection heuristics; they are not independent
        evidence if the same photometry contributed to both.
\end{itemize}

\chapter{Reproducibility and outputs}
\section{Recommended run record}
A reproducible analysis should record the target, software version, Python
version, selected interpreter, input FITS files or MAST query, sectors, cadence,
quality mask, all preprocessing parameters, period-search method and range,
stellar catalogue references, inference backend, GP kernel, random seed if set,
chain/draw/tune settings, and any warnings. The pipeline stores core settings in
results metadata and retains catalogue provenance where available.

\section{Output products}
The \code{save()} method writes results under the configured output directory,
usually in a target-specific directory such as \code{output/TIC <id>/}. Depending
on the run, products include:
\begin{itemize}
  \item JSON metadata and derived parameters;
  \item CSV tables of planet parameters and diagnostics;
  \item NetCDF ArviZ posterior data;
  \item PNG figures for each analysis stage;
  \item an HTML/PDF-style report when report generation is enabled.
\end{itemize}
The posterior NetCDF file is the authoritative record for posterior inference;
plots and summary tables are projections of that object and should not be used
as substitutes when reanalysis is required.

\section{Implementation map}
\begin{longtable}{p{0.35\textwidth}p{0.52\textwidth}}
\caption{Main implementation components.}\label{tab:files}\\
\toprule
Component & Responsibility\\
\midrule
\code{analysis/session.py} & Public staged API and complete-run orchestration\\
\code{config.py}, \code{constants.py} & Configuration precedence, validation, defaults\\
\code{data/preprocess.py} & NaN removal, protected sigma clipping, stitching, flattening\\
\code{analysis/stages/period.py} & Archive refinement and detection orchestration\\
\code{transit/detection.py} & TLS/BLS selection and iterative multi-planet search\\
\code{catalogs/stellar.py} & Catalogue merging and stellar-density calculation\\
\code{inference/priors.py} & PyMC prior definitions and parameter transforms\\
\code{inference/bayesian.py} & Keplerian transit model, GP likelihood, NUTS sampling\\
\code{inference/gp.py} & SHO or Matérn-3/2 celerite2 GP construction\\
\code{transit/parameters.py} & Posterior summaries and physical parameter derivation\\
\code{inference/diagnostics.py} & \(\hat R\), ESS, divergences, convergence flag\\
\bottomrule
\end{longtable}

\chapter{Mathematical assumptions checklist}
Before accepting a result, the analyst should explicitly answer:
\begin{enumerate}
  \item Is the input flux normalized and are TESS quality flags appropriate?
  \item Are enough independent transit events observed to distinguish \Porb from
        an integer harmonic?
  \item Was the transit mask wide enough to protect the signal but not so wide
        that most out-of-transit information was removed?
  \item Is the adopted stellar radius, mass, and density measured or fallback-
        estimated, and are their uncertainties realistic?
  \item Does a circular orbit provide a defensible approximation?
  \item Does the GP timescale separate stellar variability from the transit
        duration, and do posterior predictive checks support that separation?
  \item Are \(\hat R\), ESS, and divergences acceptable for all scientifically
        important parameters?
  \item Are the radius, temperature, and model-comparison summaries reported
        with credible intervals and provenance?
  \item Have false positives, dilution, contamination, eclipsing-binary
        scenarios, and odd--even transit differences been investigated?
\end{enumerate}

\chapter*{References and software context}
\addcontentsline{toc}{chapter}{References and software context}
\begin{description}[leftmargin=3.5cm,style=nextline]
  \item[Kipping (2013)] Efficient, uninformative sampling of limb-darkening
  coefficients in two-parameter laws, \emph{MNRAS}, 435, 2152.
  \item[Mandel and Agol (2002)] Analytic light curves for planetary transit
  searches, \emph{ApJL}, 580, L171.
  \item[Hippke et al. (2019)] Transit Least Squares, \emph{Astronomy and
  Astrophysics}, 623, A39.
  \item[Polanski et al. (2023)] TESS-Keck Survey transit-fitting methodology,
  as followed by the Bayesian implementation where applicable.
  \item[Software] The implementation uses Python, NumPy, SciPy, pandas,
  lightkurve, Astropy, astroquery, transitleastsquares, batman-package,
  PyMC, exoplanet, celerite2, ArviZ, and Matplotlib. Exact versions should be
  recorded from the active environment for each scientific run.
\end{description}

\appendix
\chapter{Parameter notation}
\begin{longtable}{p{0.18\textwidth}p{0.28\textwidth}p{0.42\textwidth}}
\caption{Notation used in this report.}\label{tab:notation}\\
\toprule
Symbol & Name & Units or meaning\\
\midrule
\(t\) & Observation time & BTJD or consistent days\\
\(F\) & Normalized flux & Dimensionless\\
\(\Porb\) & Orbital period & days\\
\(\tzero\) & Transit mid-time & BTJD\\
\(r\) & Radius ratio & \(\Rp/\Rstar\)\\
\(b\) & Impact parameter & Stellar-radius units\\
\(\Tfourteen\) & First-to-fourth contact duration & days or hours\\
\(a/\Rstar\) & Scaled semi-major axis & Dimensionless\\
\(\rhoStar\) & Stellar density & \(\mathrm{g\,cm^{-3}}\)\\
\(q_1,q_2\) & Kipping limb-darkening variables & Unit interval\\
\(\sigma_{\mathrm{GP}}\) & GP amplitude & Relative-flux units\\
\(\rho_{\mathrm{GP}}\) & GP correlation timescale & days\\
\(s\) & White-noise jitter & Relative-flux units\\
\bottomrule
\end{longtable}

\chapter{Minimal reproducibility template}
\begin{lstlisting}
from tess_pipeline import TESSAnalysis

analysis = TESSAnalysis(
    "TIC 307210830",
    inference=True,
    cadence=120,
    sectors="all",
    search_method="tls",
    inference_backend="exoplanet",
    chains=4,
    draws=1500,
    tune=1000,
    target_accept=0.95,
    gp_kernel="SHO",
)
results = analysis.run()
analysis.save("output")
\end{lstlisting}
For a controlled reanalysis, replace the target download with explicit FITS
paths, preserve the exact configuration, and archive the generated NetCDF
posterior together with catalogue query results and software-environment
metadata.

\end{document}
